% Generated by mceg equations; edit the contract, then regenerate.
% Requires amsmath. Authored math is retained; TeX is not executed here.
% Definitions and hashes do not establish scientific validity.

% equation-sha256: 87d0ea395389bae387559b03d95f6c5e93c213de0443cb20a01caf527077573a
\paragraph{Equation density\_dependent\_charge\_root (revision 1)}
Density-dependent orbital rates are evaluated inside the charge-conserving root and its implicit derivative.
\begin{equation}
\begin{aligned}
P_{ij}=R_{ij}+n_e C^{e}_{ij}+n_p C^{p}_{ij}(T,n_e),\quad n_e=n_p=n_{\rm H}e^z,\\
M_{ij}=\delta_{ij}\!\left[\sum_kP_{ik}+\Gamma_i+n_eq_i^{\rm ion}\right]-P_{ji},\\
M\boldsymbol b=n_e\boldsymbol\alpha+n_e^2\boldsymbol q^{(3)},\quad n_i=n_e b_i,\\
F(z;\theta)=z+\log\!\left(1+\sum_i b_i(n_{\rm H}e^z;\theta)\right)=0,\\
\frac{\partial z}{\partial\theta}=-\frac{F_\theta}{F_z}
\end{aligned}
\label{eq:escape.cloudy_5f_investigation:density_5f_dependent_5f_charge_5f_root:1}
\end{equation}
\noindent\textbf{Scope:} Prescribed local radiation and positive photoionization; the coupled residual uses the same rate callback.

\begin{description}
\item[$T$] Common electron and proton kinetic temperature (K; ID: temperature).
\item[$n_e$] Electron density, equal to proton density in pure hydrogen (cm\textasciicircum{}-3; ID: electron\_density).
\item[$n_{\rm H}$] Total hydrogen-nucleus density (cm\textasciicircum{}-3; ID: density).
\item[$z$] Log ionized fraction (1; ID: log\_fraction).
\item[$R,C,\Gamma,\alpha,q^{\rm ion},q^{(3)}$] Radiative, binary and three-body rate coefficients (s\textasciicircum{}-1, cm\textasciicircum{}3 s\textasciicircum{}-1, cm\textasciicircum{}6 s\textasciicircum{}-1; ID: rates).
\item[$b_i$] Bound population divided by proton density (1; ID: ratios).
\item[$\theta$] External density, temperature and radiation inputs (various; ID: parameters).
\end{description}

\noindent\textbf{Assumptions:}
\begin{itemize}
\item Pure H, ne=np; the bound-state loss matrix is nonsingular.
\item The bracket is -80<=z<=0. The locally positive F\_z is checked; derivatives describe a converged root, not the bisection iterations.
\end{itemize}

% equation-sha256: 3d1b4edb827c916a9793cf31284e44806dbceaaa109656d55d4ffeec04d80e21
\paragraph{Equation psm20\_cross\_section (revision 1)}
PSM20 final-state-resolved impact-parameter cross section with both cutoffs integrated continuously.
\begin{equation}
\begin{aligned}
u(E)=\sqrt{\frac{E}{\mathrm{Ry}}\frac{m_e}{\mu}},\quad
D=\frac{3n^2\ell(n^2-\ell^2)}{2\ell+1},\quad R_1=a_0\sqrt{\frac{4D}{u^2}},\\
R_D=\sqrt{\frac{k_{\rm B}T}{4\pi n_e e^2}},\quad
R_\tau=0.72\tau\frac{a_0}{t_0}u,\quad R_c=\min(R_D,R_\tau),\\
\sigma(E)=\frac{\pi R_1^2}{4}\begin{cases}
\frac23+2\log(R_c/R_1),&R_c\ge R_1,\\
\frac23(R_c/R_1)^3,&R_c<R_1,
\end{cases}\\
q_{n\ell\to n,\ell-1}=\sqrt{\frac{8}{\pi\mu}}(k_{\rm B}T)^{-3/2}
\int_0^\infty E\sigma(E)e^{-E/(k_{\rm B}T)}\,dE
\end{aligned}
\label{eq:escape.cloudy_5f_investigation:psm20_5f_cross_5f_section:1}
\end{equation}
\noindent\textbf{Scope:} Proton-driven n=2,3 orbital mixing; same-n states are degenerate in the cross section, with actual energy differences in the inverse rate.

\begin{description}
\item[$T$] Common electron and proton kinetic temperature (K; ID: temperature).
\item[$n_e$] Electron density, equal to proton density in pure hydrogen (cm\textasciicircum{}-3; ID: electron\_density).
\item[$n$] Principal quantum number (1; ID: principal).
\item[$\ell$] Departing orbital quantum number (1; ID: orbital).
\item[$\ell-1$] Final orbital quantum number (1; ID: lower\_orbital).
\item[$\tau$] Departing state spontaneous lifetime (s; ID: lifetime).
\item[$\mu$] Neutral-protium/proton reduced mass (g; ID: mass).
\item[$E$] Collision kinetic energy (erg; ID: energy).
\item[$q$] Maxwellian binary rate coefficient (cm\textasciicircum{}3 s\textasciicircum{}-1; ID: rate).
\item[$R_1,R_D,R_\tau,R_c$] Matching, electron-Debye, lifetime and selected cutoff radii (cm; ID: radii).
\item[$a_0,t_0,\mathrm{Ry}$] Bohr length, time and Rydberg energy (cm, s, erg; ID: atomic\_units).
\end{description}

\noindent\textbf{Assumptions:}
\begin{itemize}
\item P1=1/4 for one resolved final state; the projectile charge is one.
\item The lifetime is computed from the sum of spontaneous transitions out of the departing state.
\item Unlike the C23 shortcut, the lifetime/Debye minimum is integrated at every density; this intentionally removes the shortcut jump.
\item No electron orbital-mixing channel or population fit is added to this chosen finite-atom closure.
\end{itemize}

% equation-sha256: 68d785bc96d6c29b220870e92084ae4be4e0d6e7e88a7feac9605aa66593cb99
\paragraph{Equation mixing\_inverse\_pair (revision 1)}
The paired orbital transitions obey the finite atom energy and statistical-weight convention.
\begin{equation}
\begin{aligned}
C^p_{lu}=\frac{g_u}{g_l}\exp\!\left[-\frac{E_u-E_l}{k_{\rm B}T}\right]C^p_{ul},\\
(u,l)\in\{(2p,2s),(3p,3s),(3d,3p)\},\qquad
P^{\rm mix}_{ij}=n_pC^p_{ij}
\end{aligned}
\label{eq:escape.cloudy_5f_investigation:mixing_5f_inverse_5f_pair:1}
\end{equation}
\noindent\textbf{Scope:} The finite atom replaces, rather than adds to, the previous 2s-2p fit. The historical six-state API keeps that fit.

\begin{description}
\item[$T$] Common electron and proton kinetic temperature (K; ID: temperature).
\item[$n_e$] Electron density, equal to proton density in pure hydrogen (cm\textasciicircum{}-3; ID: electron\_density).
\item[$C^p_{ij}$] Proton collision coefficient (cm\textasciicircum{}3 s\textasciicircum{}-1; ID: coefficient).
\item[$g_i$] Finite-state statistical weight (1; ID: weight).
\item[$E_i$] Finite-state excitation energy (erg; ID: excitation).
\end{description}

\noindent\textbf{Assumptions:}
\begin{itemize}
\item Pure H makes np=ne, so the paired proton rates may join the effective electron collision matrix.
\item Collapsed n=4,5 shells retain statistical orbital populations.
\end{itemize}

% equation-sha256: 31370290b36e7174c8c8b0a2348b4369761e915a926df273098add77780dcde4
\paragraph{Equation hydrogenic\_bound\_free (revision 1)}
Coulomb-wave bound-free dipoles resolve orbital states; collapsed shells average their cross sections statistically.
\begin{equation}
\begin{aligned}
t^2=\frac{h\nu}{\chi_{n\ell}}-1,\quad \kappa=t/n,\\
\sigma_{n\ell}=\frac{4\pi\alpha a_0^2n^2}{3}
\sum_{\ell'=\ell\pm1\ge0}\frac{\max(\ell,\ell')}{2\ell+1}(1+t^2)
\mathcal D_{n\ell,\ell'}^2,\\
\mathcal D_{n\ell,\ell'}^2=\frac{(n+\ell)!}{(n-\ell-1)!}(2n)^{2(\ell-n)}
\prod_{s=1}^{\ell'}(1+s^2\kappa^2)[G_\ell^{\pm}]^2,\\
\overline\sigma_n=\sum_{\ell=0}^{n-1}\frac{2\ell+1}{n^2}\sigma_{n\ell}
\end{aligned}
\label{eq:escape.cloudy_5f_investigation:hydrogenic_5f_bound_5f_free:1}
\end{equation}
\noindent\textbf{Scope:} 1<=n<=5; sigma=0 below the actual finite-state threshold. The radial functions are defined by the following recurrence.

\begin{description}
\item[$n$] Principal quantum number (1; ID: principal).
\item[$\ell$] Bound-state orbital quantum number (1; ID: orbital).
\item[$h\nu/\chi_{n\ell}$] Photon energy divided by the level threshold (1; ID: relative\_energy).
\item[$\nu$] Photon frequency (Hz; ID: frequency).
\item[$\sigma_{n\ell}$] Orbital-resolved photoionization cross section (cm\textasciicircum{}2; ID: cross\_section).
\item[$\mathcal D,G^\pm$] Dimensionless Coulomb radial factors (1; ID: radial).
\item[$\chi_{n\ell}$] Finite-level ionization threshold (erg; ID: threshold).
\end{description}

\noindent\textbf{Assumptions:}
\begin{itemize}
\item Resolved states are 1s,2s,2p,3s,3p,3d; n=4,5 are collapsed.
\item Hydrogenic Coulomb dipoles are evaluated at the selected level threshold.
\item No fitted correction is applied to n2/n3 populations.
\end{itemize}

% equation-sha256: 9861b547e50970f84ebd2d0f5827714e37bcb847a1ad02bfc057b015cda2eddf
\paragraph{Equation brocklehurst\_recurrence (revision 1)}
Finite downward recurrence used to evaluate the hydrogenic radial factors without population fitting.
\begin{equation}
\begin{aligned}
G_0^*=\sqrt{\frac\pi2}\frac{8n(4n)^n}{(2n-1)!}
\frac{e^{-2n\arctan(t)/t}}{\sqrt{1-e^{-2\pi/\kappa}}(1+t^2)^{n+2}},\\
G^+_{n-1}=G_0^*,\quad G^+_{n-2}=n(2n-1)(1+t^2)G_0^*,\\
G^-_{n-1}=\frac{1+t^2}{2n}G_0^*,\quad
G^-_{n-2}=(2n-1)[4+(n-1)(1+t^2)]G^-_{n-1},\\
A_q^+=4n^2-4(q+2)^2+(q+2)(2q+3)(1+t^2),\\
B_q^+=4n^2[n^2-(q+2)^2][1+(q+3)^2\kappa^2],\\
A_q^-=4n^2-4(q+1)^2+(q+1)(2q+3)(1+t^2),\\
B_q^-=4n^2[n^2-(q+2)^2][1+(q+1)^2\kappa^2],\\
G_q^\pm=A_q^\pm G_{q+1}^\pm-B_q^\pm G_{q+2}^\pm
\end{aligned}
\label{eq:escape.cloudy_5f_investigation:brocklehurst_5f_recurrence:1}
\end{equation}
\noindent\textbf{Scope:} The plus recurrence runs q=n-3,...,0 and the minus recurrence q=n-3,...,1. Use initial values only for existing indices.

\begin{description}
\item[$n$] Principal quantum number (1; ID: principal).
\item[$\ell$] Bound-state orbital quantum number (1; ID: orbital).
\item[$1+t^2$] Photon energy divided by threshold (1; ID: relative\_energy).
\item[$G_q^\pm$] Radial recurrence quantities (1; ID: radial).
\item[$\kappa=t/n$] Dimensionless continuum momentum (1; ID: continuum).
\end{description}

\noindent\textbf{Assumptions:}
\begin{itemize}
\item At threshold arctan(t)/t approaches one and exp(-2pi/kappa) approaches zero.
\item Small-t series avoids a singular intermediate and preserves differentiation.
\item Ordinary float64 is restricted to n<=5; high-n overflow behavior is not inferred.
\end{itemize}

% equation-sha256: 9e5010ac050dd8f0a805455aa753cc2eea3484186b90bea6cf1d0677c49d9024
\paragraph{Equation paired\_continuum\_kernels (revision 1)}
One photoionization cross section supplies its inverse Milne rate and continuum absorption/emission.
\begin{equation}
\begin{aligned}
\eta_\nu=\frac{c^2J_\nu}{2h\nu^3},\quad
 a_i=\frac{g_i}{2}\left(\frac{h^2}{2\pi m_e k_{\rm B}T}\right)^{3/2}
 e^{-(h\nu-\chi_i)/(k_{\rm B}T)},\\
\Gamma_i=\frac{8\pi}{c^2}\int\nu^2\sigma_i\eta_\nu\,d\nu,\quad
\alpha_i=\frac{8\pi}{c^2}\int\nu^2\sigma_i a_i(1+\eta_\nu)\,d\nu,\\
\chi_\nu^{\rm bf}=\sum_i\sigma_i(n_i-n_en_pa_i),\quad
j_\nu^{\rm bf}=\frac{2h\nu^3}{c^2}n_en_p\sum_i\sigma_i a_i
\end{aligned}
\label{eq:escape.cloudy_5f_investigation:paired_5f_continuum_5f_kernels:1}
\end{equation}
\noindent\textbf{Scope:} Included bound-free processes only; other free-free and line kernels are unchanged and separately recorded in the earlier coupled formulation.

\begin{description}
\item[$T$] Common electron and proton kinetic temperature (K; ID: temperature).
\item[$n_e$] Electron density, equal to proton density in pure hydrogen (cm\textasciicircum{}-3; ID: electron\_density).
\item[$\nu$] Frequency (Hz; ID: frequency).
\item[$\eta_\nu$] Photon occupation number (1; ID: occupation).
\item[$d\nu$] Frequency quadrature measure (Hz; ID: weight).
\item[$n_i$] Bound population (cm\textasciicircum{}-3; ID: population).
\item[$\sigma_i$] The same orbital-resolved or collapsed cross section in every kernel (cm\textasciicircum{}2; ID: cross\_section).
\item[$\Gamma_i,\alpha_i$] Photoionization and radiative recombination coefficients (s\textasciicircum{}-1, cm\textasciicircum{}3 s\textasciicircum{}-1; ID: rates).
\item[$\chi_\nu^{\rm bf},j_\nu^{\rm bf}$] Bound-free extinction and emissivity (cm\textasciicircum{}-1, erg s\textasciicircum{}-1 cm\textasciicircum{}-3 sr\textasciicircum{}-1 Hz\textasciicircum{}-1; ID: transfer).
\end{description}

\noindent\textbf{Assumptions:}
\begin{itemize}
\item Maxwellian electrons at the gas temperature; stimulated recombination is retained.
\item The same selected quadrature and cross section must be used in local rates and coupled continuum transfer.
\end{itemize}
